\chapter{Dracunculiasis (Guinea Worm Disease)}
\index{Dracunculiasis (Guinea Worm Disease)}

\section*{Definition and Overview}
Dracunculiasis, commonly called guinea worm disease, is a parasitic infection caused by the nematode Dracunculus medinensis. It is acquired by drinking water containing tiny water fleas (copepods) that carry the immature worms. Over about a year, the worms grow inside the body, and the female worm eventually emerges through a painful blister, usually on the leg or foot. Guinea worm disease causes severe pain and disability and is a disease of poverty and unsafe water. It is close to being eradicated worldwide through water filtration, health education, and treatment of cases. This chapter explains the infection, its effects, and its prevention.

\section*{Epidemiology}
Guinea worm disease has declined dramatically through the global eradication campaign. In the 1980s, around 3.5 million cases occurred each year in 20 countries; by recent years, only a handful of cases were reported, mainly in a few countries in Africa. The disease occurs in poor, rural communities that rely on unprotected water sources. It mainly affects adults who work in farming, and each infection causes about a month of disability. some countries have no guinea worm disease, and it is only relevant to travellers from affected regions.

\section*{Aetiology and Pathophysiology}
\begin{itemize}
    \item \textbf{Parasite:} Dracunculus medinensis, a long thin worm, is the cause
    \item \textbf{Transmission:} drinking water containing copepods infected with larvae
    \item \textbf{Development:} the larvae mature in the abdomen over about a year
    \item \textbf{Emergence:} the female worm travels to the skin and emerges through a blister, releasing larvae when the wound touches water
    \item \textbf{Secondary infection:} the emerging wound often becomes infected with bacteria
    \item \textbf{Eradication:} the disease is preventable and has no animal reservoir, making eradication achievable
\end{itemize}

\section*{Clinical Features}
\begin{itemize}
    \item No symptoms for about a year after infection
    \item A burning sensation and a painful blister, usually on the lower leg or foot
    \item The blister ruptures when in contact with water
    \item The worm emerges slowly over weeks
    \item Fever, nausea, vomiting, and dizziness as the worm emerges
    \item Severe pain, swelling, and secondary bacterial infection of the wound
    \item Joint stiffness and disability lasting weeks to months
\end{itemize}

\section*{Differential Diagnosis}
\begin{itemize}
    \item Other causes of skin blisters and ulcers, including boils and tropical ulcers
    \item Other parasitic infections, including filariasis and loiasis
    \item Cellulitis and other skin infections
    \item Joint and muscle conditions causing pain and disability
\end{itemize}

\section*{When to Seek Medical Care}
Travellers who develop a painful blister after returning from affected regions should see a doctor and mention their travel history. Anyone with a suspected guinea worm infection needs medical care for the worm's removal and treatment of complications.

\textbf{Contact your local emergency services for:}
\begin{itemize}
    \item Signs of severe infection: spreading redness, high fever, or feeling very unwell
    \item Inability to walk or severe swelling
    \item Any medical emergency
\end{itemize}

\section*{Diagnosis and Investigations}
\begin{itemize}
    \item \textbf{Clinical diagnosis:} the emerging worm and the history of exposure are characteristic
    \item \textbf{Examination:} identification of the worm in the blister or wound
    \item \textbf{Blood tests:} may show elevated eosinophils
    \item \textbf{Imaging:} X-rays can show calcified dead worms
    \item \textbf{Travel history:} essential for recognition outside endemic areas
\end{itemize}

\section*{Management}
\begin{itemize}
    \item \textbf{Slow worm removal:} the worm is drawn out gradually over days to weeks, wrapped around a stick
    \item \textbf{Never break the worm:} breaking it causes severe inflammation
    \item \textbf{Wound care:} cleaning and dressings prevent secondary infection
    \item \textbf{Pain relief:} analgesics for the pain
    \item \textbf{Antibiotics:} for secondary bacterial infection
    \item \textbf{Keeping the wound out of water:} prevents the worm releasing larvae and spreading the disease
    \item \textbf{Prevention of spread:} education and water safety in affected communities
\end{itemize}

\section*{Complications}
\begin{itemize}
    \item Secondary bacterial infection of the wound, including tetanus
    \item Abscess formation where the worm is broken
    \item Joint inflammation and stiffness
    \item Prolonged disability and lost work
    \item Tetanus, which is preventable by vaccination
\end{itemize}

\section*{Prognosis and Outcomes}
Guinea worm disease is not usually fatal, but it causes significant pain and disability, typically for a month or more. With careful wound care and gradual worm removal, recovery is complete, although secondary infection can complicate healing. Because the disease has no animal reservoir, it is fully preventable, and the global eradication programme has brought it to the brink of elimination, with the ultimate goal of making guinea worm disease the second human disease ever eradicated.

\section*{Prevention and Self-Care}
\begin{itemize}
    \item Drink only filtered or treated water in affected regions
    \item Filter drinking water through fine cloth to remove the water fleas
    \item Keep anyone with an emerging worm away from water sources
    \item Support safe water supplies and health education in affected communities
    \item Keep tetanus vaccination up to date for travellers
    \item Report any suspected case to health authorities
    \item Follow the medical team's advice for wound care
\end{itemize}

\section*{Sources and Further Reading}
The following reputable sources provide further information for healthcare students and patients seeking reliable, current advice about guinea worm disease.

\begin{itemize}
    \item World Health Organization. \textit{Dracunculiasis (guinea worm disease) fact sheet.} https://www.who.int/
    \item The Carter Center. \textit{Guinea Worm Eradication Program.} https://www.cartercenter.org/
    \item Centers for Disease Control and Prevention. \textit{Guinea worm disease.} https://www.cdc.gov/
    \item MSD Manual. \textit{Dracunculiasis: professional overview.} https://www.msdmanuals.com/
    \item MedlinePlus. \textit{Guinea worm infection.} https://medlineplus.gov/
    \item NSW Health. \textit{Travellers' health information.} https://www.health.nsw.gov.au/
\end{itemize}
